\scanpage{481}

For this second point, we have an immediate answer (assuming the
$\unsure{\mathcal{N}_{g,\nu}|N_{g,\nu}}$ already known). Let \unsure{a|a} $(\pi,a)$ be of type $g,\nu$,
with \deleted{\unsure{when}} a \inserted{partial} lifting of $\boldsymbol{\Gamma}_a$ \unsure{to|a}
$\Gamma'_{\unsure{a|\cdot}}\hookrightarrow \unsure{\mathcal{N}_a|N_a}$.
It will be admissible \unsure{iff} we can find
a $(\pi',a')$ of type $(g,\nu')$, \inserted{and} an
embedding of $\pi'$ as a \unsure{subgroup|subgroup} of finite index in $\widehat{\pi}_{03}$ compatible with $a'$,
\unsure{and} an open subgroup $\unsure{\Gamma}$ \inserted{$\Gamma''$} of $\boldsymbol{\Gamma}_{03}$, and
an \inserted{admissible} actual \unsure{action} of $\Gamma''$ \deleted{\unsure{admissible}}
on $\widehat{\pi}_{03}$ which
lifts the given outer action, and which
leaves $\pi'$ invariant, so that up to isomorphism,
$(\pi,a)$ \inserted{and the germ of action of $\Gamma'$ on it}
is deduced from $(\pi',a')$
by the action of $\Gamma''$ on it, by ``plugging''
a batch of holes \inserted{(and \unsure{forgetting} an actual action in favor
of the \unsure{outer} action)}. One needs
in this approach simply to know how to
specify \inserted{\unsure{algebraically}} what is meant by
an ``admissible'' action of \inserted{$\Gamma''\subset$} $\boldsymbol{\Gamma}_{0,3}$ on $\widehat{\pi}_{0,3}$
--- it being understood that this must correspond
to the points of $\unsure{U_{0,3}}$ rational over \deleted{$\unsure{\mathbb{Q}}$} the number
field defined by $\Gamma''$. One can conjecture

\margin{480}

\scanpage{482}
\mspage{244}

this one, by ``plugging the holes''.

It seems that on the same
principle one can give a description of the
$\unsure{\mathcal{N}_{g,\nu}|N_{g,\nu}}$ (hence of the $\boldsymbol{\Gamma}_{g,\nu}$)
in terms of
$\boldsymbol{\Gamma}_{0,3}$. Using the transition homomorphisms
$\unsure{\widehat{\widehat{\mathcal{T}}}_{g,\nu}}
\longrightarrow
\widehat{\widehat{\mathcal{T}}}_{g,\nu'}$,
it suffices to
do so, when $g$ is fixed, for large
$\nu$ --- large enough that there
exists an \inserted{algebraic} curve of type $g,\nu$ \inserted{over $\mathbb{Q}$}
which is an étale covering of $U_{0,3}$ over $\mathbb{Q}$.

We therefore consider a \unsure{corresponding} embedding
of $\pi_{g,\nu}$ into $\pi_{0,3}$, and
decree that if \deleted{$\unsure{\boldsymbol{\Gamma}_{03}|\mathcal{T}_{03}}$} we can make
$\mathrm{ext}^{t}$ $\boldsymbol{\Gamma}_{0,3}$ operate on $\pi_{g,\nu}$, so
that $\pi_{g,\nu}\longrightarrow\pi_{0,3}$ \unsure{commutes}
with these outer actions, \uline{then} the subgroup
of $\widehat{\widehat{\mathcal{T}}}_{g,\nu}$ generated by
$\widehat{\mathcal{T}}_{g,\nu}$ and
$\boldsymbol{\Gamma}_{0,3}$ is \deleted{$\unsure{\widehat{\mathcal{T}}}$}
$\unsure{\mathcal{N}_{g,\nu}|N_{g,\nu}}$.

\margin{481}

\scanpage{483}
\mspage{245}
34. Heuristic profinite description of the category
of algebraic curves \unsure{defined} over \deleted{\unsure{$\mathbb Q$}} finite
subextensions $K$ of $\overline{\mathbb Q}_0/\mathbb Q$ (i.e. of $\mathbb C/\mathbb Q$).

We restrict ourselves to \inserted{geometrically} connected curves \inserted{(for convenience)}, anabelian
(out of provisional necessity \unsure{!}, cf. below on \unsure{the way} to
get rid of this restriction). The datum of
a universal covering of \deleted{a} $U_{\overline{\mathbb Q}_0}$ defines an
extension structure
\[
(1)\qquad 1\longrightarrow\pi\longrightarrow\mathcal E\longrightarrow\Gamma^{+}\longrightarrow1\qquad\text{or}
\]
(2) or again we have a homomorphism $\mathcal E\longrightarrow\boldsymbol{\Gamma}$
(image $\Gamma^{+}$ open in $\boldsymbol{\Gamma}$,
kernel called $\pi$),
where $\Gamma^{+}\subset\boldsymbol{\Gamma}_{\mathbb Q}\simeq\boldsymbol{\Gamma}_{0,3}$ is the open \unsure{subgroup|subgroup}
corresponding to $K$. Giving oneself such an
extension (assuming $\unsure{\mathrm{Centre}}(\pi)=1$) amounts to
giving oneself an outer action of $\Gamma'$ on $\pi$.
A first question: must one put the
loop structure of $\pi$ in the \deleted{\illegible{1}} data \unsure{$\mathcal E$} of
the object (1) (or \unsure{(2)|(1')}) supposed to describe $U\unsure{\mathcal E}/K$?

\uline{Conjecture} It is not worth the trouble -- the loop structure
of $\pi$ is the only loop structure \unsure{invariant}
under the outer action of $\Gamma^{+}$ (or of $\Gamma^{\natural}$). The
loop subgroups $L_i$ are the maximal
subgroups\inserted{in $\pi$}, such that $\mathrm{Norm}_{\mathcal E}(L_i)\longrightarrow\Gamma^{+}$ has
as image an open subgroup of $\unsure{\Gamma^{+}|\Gamma}$.
\margin{\unsure{and} every subgroup $\ne1$ of $\pi$ which has this property is contained in a unique $L_i\ldots$}
I
also presume that every homomorphism \deleted{\unsure{of}}\inserted{between} extensions
$\mathcal E$ of $\Gamma^{+}$ by a $\pi$, $\mathcal E'$ of $\Gamma$ by a $\pi'$
($\mathcal E,\mathcal E'$ coming from algebraic curves $U,U'$) and such
that the image of $\mathcal E$ in $\mathcal E'$ is open, respects
\margin{482}

\scanpage{484}
necessarily the loop structure, and in fact
comes from a morphism (unique, as we know) of algebraic
curves.

At any rate, if one admits the description of the
loop subgroups, it \unsure{will be} clear that the image under
$\unsure{u}$ of an $L_i$ will be, either $(1)$, or contained
in a unique $L'_j\ldots$

This would mean that the functor from algebraic
curves over $K$\inserted{(morphisms the \unsure{dominant} morphisms)}, to profinite outer groups
on which $\Gamma$ operates, would be fully faithful.
The functor ``extension of the base field'' from
$K$ to $K'$ corresponds to the functor: restriction
of an \unsure{outer|outer} group (or of an extension
structure) from $\Gamma$ to $\Gamma'$ (open subgroup).
The finite \inserted{\unsure{étale}} coverings of $U$ correspond to
$\mathcal E$-sets: [$\widetilde U$ \unsure{being chosen}]\ldots

To describe the essential image of this
functor, we are not reduced to conjectures:
we start from the \inserted{canonical} extension $\mathcal E_{03}$ of $\boldsymbol{\Gamma}_{03}=\boldsymbol{\Gamma}$
by $\widehat{\pi}_{03}$, we take an \inserted{open} \unsure{subgroup|subgroup} $\mathcal E'$
of $\mathcal E_{03}$, with image $\Gamma\subset\boldsymbol{\Gamma}_{03}$, kernel
$\pi'\subset\widehat{\pi}_{0,3}$, and in the canonical
loop structure of $\pi'$\inserted{of genus $g$} (induced by that of $\widehat{\pi}_{03}$)
we take an $\unsure{I|\overline I}\subset I(\pi')$\inserted{such that $2g+\unsure{\mathrm{Card}(I)}\geq3$}, \inserted{stable under $\Gamma$, and} we ``plug the holes''
at $I(\pi')\setminus I$. Likewise, to describe when
\margin{483}

\scanpage{485}
\mspage{\unsure{246|245}}
an actual action, lifting an admissible outer
action, is itself admissible -- \deleted{\unsure{can}}
i.e. can be obtained from an algebraic \deleted{\illegible{1}}\inserted{curve}
$U$, \uline{pointed} by a point rational over $K$.
This does not \uline{mean} for all that, that if
$U$ is given by an outer action of $\Gamma$ on a $\pi$,
that the conjugacy classes of admissible liftings
of this action indeed all \unsure{come} from
points of $U$ rational/$K$. But we see at once that
this will be the case, as soon as one admits
full faithfulness for \uline{isomorphisms}

\uline{N.B.} Even for the automorphisms of $U_{0,3}$
$=\mathbb P^1_{\mathbb Q}\setminus\{0,1,\infty\}$, this full faithfulness is not
at all clear. One would have to prove that \deleted{\illegible{1}}
the commutant of $\boldsymbol{\Gamma}_{0,3}$ in $\widehat{\widehat{\mathcal T}}_{0,3}$ is reduced
to $\unsure{\mathfrak S_3|\mathcal S_3}$. The situation is less hopeless,
when we give ourselves, along with the structure of
a \inserted{profinite} outer group with operators $\Gamma$, an arithmetization
of $\pi$ (which supposes that we have \unsure{made explicit}
a loop structure) invariant under \deleted{\illegible{1}} $\Gamma$
-- \unsure{thus} in $\widehat{\widehat{\mathcal T}}(\pi)$ we have at our disposal an $\unsure{N(\pi)|\mathcal N(\pi)}$,
and $\Gamma\hookrightarrow\unsure{N(\pi)|\mathcal N(\pi)}$. From this point of view, \inserted{for} a
homomorphism \inserted{$\pi'\longrightarrow\pi$} of outer $\Gamma$-groups (arithmetized),
\unsure{it} is understood that it is compatible with
the arithmetizations, i.e. if $\pi''$ is the image of $\pi'$ in $\pi$,
\margin{484}

\scanpage{486}

we see that the arithmetization of $\pi''$ \unsure{deduced} from that of $\pi'$ by ``passage to the quotient'', be the one \unsure{induced} by $\pi$. \deleted{\unsure{The interest}} In particular, \inserted{for} the automorphisms\inserted{(outer)} of $\pi$ it is understood that \deleted{\illegible{1}} not only do they commute with $\Gamma$, but also that they are in $N(\pi)$. \inserted{(which implies that they are in \unsure{$\widehat{\mathcal T}(\pi)$}, for the centralizer in $\boldsymbol{\Gamma}_\pi$ of \unsure{every} open subgroup of $\boldsymbol{\Gamma}_\pi$ is $(1)$ !)} In the case of $U_{0,3}$, we have $N_{0,3}\simeq\unsure{\mathfrak S_3}\times\boldsymbol{\Gamma}_{0,3}$, and we know that the center of $\boldsymbol{\Gamma}_{03}\simeq\boldsymbol{\Gamma}$ is trivial -- thus we do indeed find that the group of automorphisms of this structure is reduced \unsure{to|:} \unsure{$\mathfrak S_3$} !

This point of view is nevertheless probably superfluous -- for we \unsure{presume} that for the given outer action of $\Gamma$, there is a unique arithmetization invariant under $\Gamma$, and that the \inserted{``admissible''} \unsure{homomorphisms} of the ``admissible'' outer $\Gamma$-groups, as they have been defined previously, \unsure{automatically} respect this arithmetization. If this were not so, one would of course have to introduce the arithmetizations into the structure.

\deleted{\unsure{Here}, \deleted{\unsure{as} $\unsure{\boldsymbol{\Gamma}_\pi|\Gamma_\pi}\unsure{=}$} \inserted{if $\Gamma=\unsure{\boldsymbol{\Gamma}_{\mathbb Q}}$,} the arithmetization of \deleted{\illegible{1}} $\pi$ \deleted{\unsure{is}} \inserted{must be} \unsure{described} in fact by a well-determined oriented prediscretification $\alpha$, \deleted{\unsure{from the fact that}\illegible{2}} \inserted{In fact, \unsure{we are going to|we for ourselves} describe an \unsure{oriented} \uline{strict} prediscretification} \deleted{\illegible{3} \unsure{base point}} well-determined of $\pi$, Let us start from an admissible lifting of the given outer operation}

\margin{485}

\scanpage{487}
\mspage{247}

The question arises of finding a more \deleted{\unsure{simple and}} \inserted{simplistic} description of the outer actions of a $\Gamma$ on a $\pi$ which are ``admissible''. \deleted{\illegible{1}} Here, we shall start from a $\pi$ for which we already \unsure{fix} a loop structure \inserted{and an arithmetization \unsure{a}} (\unsure{invariant|invariant} under $\Gamma$), so \deleted{\illegible{2}} that $\Gamma\subset N(\pi)$, $\Gamma\longrightarrow\boldsymbol{\Gamma}_\pi$ \unsure{injective} with open image, \deleted{\illegible{1}} i.e. $\Gamma$ corresponds to a partial splitting (or germ of splitting) of \deleted{\illegible{1}}
\[
1\longrightarrow\widehat{\mathcal T}(\pi)\longrightarrow N(\pi)\longrightarrow\boldsymbol{\Gamma}_\pi\longrightarrow1
\]

\unsure{Let} $(g,\nu)$ be the type of $\pi$; \deleted{\illegible{2}} if $g\geq1$ the simplest criterion, it is that the ``\inserted{$\Gamma^{\natural}$-}points'' \unsure{of|of the} $\pi'$ deduced from $\pi$ by plugging all the holes, be distinct. \deleted{\unsure{Except}} (N.B. if $g=1$, \deleted{\unsure{this}} \deleted{\unsure{plugging}} \inserted{by plugging} all the holes we land in an abelian case -- but that does not matter). If $g=0$, \deleted{the criterion is that} there is no condition for $\unsure{\alpha|\nu}=3$, and if $\nu>3$, \unsure{let us set|setting} aside a \deleted{\unsure{a}} subset $I'\subset I(\pi)$ with $\mathrm{card}\,I'=3$, the condition is that on plugging the holes \unsure{of|in} $\unsure{I\setminus I'}$, the \inserted{$\Gamma^{\natural}$-}points \unsure{of the|of} $\pi$ deduced from the points of $\unsure{I\setminus I'}$ are distinct.

One will note that in this conjectural perspective, in the limiting cases $(g,0)$ $(g\geq2)$, $(1,1)$, $(0,3)$, we impose no condition \inserted{a priori} on the \unsure{liftings}.

\margin{486}

\scanpage{488}

This is perhaps very \unsure{brutal} -- and it could be that we find liftings which do \unsure{not} \unsure{correspond} to an algebraic curve -- i.e. which in fact are \unsure{not} admissible, even if they appear to be.

To understand something of this, it seems to me that we must return to an interpretation of the \inserted{germs of} splittings called ``admissible'' of an extension such as
\diagram{\begin{tikzcd}[ampersand replacement=\&]
1 \arrow[r] \& \widehat{\mathcal T}_{g,\nu} \arrow[r] \& N_{g,\nu} \arrow[r] \& \boldsymbol{\Gamma}_{g,\nu} \arrow[r] \& 1 \\
\& \& \& \boldsymbol{\Gamma}_{\mathbb Q}\simeq\boldsymbol{\Gamma}_{0,3} \arrow[u,"\unsure{\wr}"'] \&
\end{tikzcd}}
as corresponding to the algebraic points of a modular \uline{variety} (\unsure{rather here}, a multiplicity) \unsure{$M_{g,\nu\,\mathbb Q}$|$M_{g,\nu,\mathbb Q}$}, of which $N_{g,\nu}$ is the arithmetic fundamental group, and $\widehat{\mathcal T}_{g,\nu}$ the geometric fundamental group.

I would like to examine this situation more closely, later on. For the moment, it seems prudent not to make \unsure{any|any} hasty conjectures for \deleted{\illegible{1}} a \uline{direct} description (i.e. \uline{not} via $U_{03}$) of ``admissible'' outer operations. We shall therefore work for the moment \unsure{with|on} this notion in the ``constructive'' form (via $U_{03}$).

\margin{487}

\scanpage{489}
\mspage{248}

If we have an \inserted{outer} $\pi$ of type \unsure{$(g\nu)$|$(g,\nu)$} with an admissible outer operation of $\Gamma\subset\boldsymbol{\Gamma}$, we expect that there will be a canonical \unsure{oriented} prediscretification on \unsure{$\pi_x$} (not \unsure{invariant|invariant} under $\Gamma$, of course \unsure{!|)} -- but such that the corresponding arithmetization is). We \unsure{shall} \inserted{even}, for an admissible \inserted{actual} action of $\Gamma^{\natural}$ on $\pi$, \inserted{(lifting the given outer action)}, define a canonical \unsure{corresponding} \inserted{oriented} discretification $\pi_0\subset\pi$ (\unsure{replaced} by a \unsure{conjugate}, when we conjugate \deleted{\unsure{the action}} the lifting $\Gamma^{\natural}\to\mathcal E$ by a $g\in\pi$) -- all these oriented discretifications (corresponding to the different ``points'' of $B_{\pi,\Gamma^{\natural}}$) define one and the same oriented \inserted{pre-}discretification -- and even one and the same \uline{strict} oriented prediscretification.

The actual action of $\Gamma^{\natural}$ on $\pi$ can be obtained (up to isomorphism) \inserted{as follows:} we \unsure{take|shall take} an \inserted{open} subgroup $\mathcal E'\subset\mathcal E_{03}$, whence extension
\[
1\to\pi'\to\mathcal E'\to\Gamma'\to1
\]
with $\pi'\subset\widehat{\pi}_{03}$, hence $\pi'=\widehat{\pi}'_0$ $(\pi'_0=\pi'\cap\unsure{\pi|\Pi}_{03})$. We have an operation of $\Gamma'$ on $I(\pi')=I(\pi'_0)$, we take it trivial (at the cost of passing to an open subgroup of $\Gamma'$), we choose $i\in I'\subset I(\pi')$, we plug the holes at $I'$, \deleted{and we look at} \inserted{whence $\pi$, an \deleted{outer action}} \inserted{oriented} discretification $\pi_0$, and an outer action of $\Gamma'$ on it, which is lifted to an actual action thanks to $i$. We do indeed find in this way a

\margin{488}

\scanpage{490}

``lattice'' \inserted{-- an oriented discretification} in $\pi$ -- I say that it does not depend on the choices that \unsure{have} been made -- in particular, that the \unsure{automorphisms|automorphisms} of $(\pi,a,\Gamma^{\natural})$ \unsure{carry} $\pi_0$ into itself. Moreover, if we have two points \inserted{$x,y$} of $B_{\pi,\Gamma^{\natural}}$, whence $\pi(x)$, $\pi(y)$, among the \unsure{classes} of paths from $x$ to $y$ (which \unsure{make|form} a \unsure{bitorsor} under $\pi(y),\pi(x)$) there are some for which \deleted{\illegible{2}} $\pi(x)\to\pi(y)$ sends $\pi_0(x)$ into $\pi_0(y)$, -- when we \unsure{restrict} \unsure{ourselves} to these, we find a subgroupoid of the fundamental groupoid $\Pi_{\unsure{1|\cdot}}(\pi,\Gamma^{\natural})$, which this time is a \uline{\unsure{connected}} loop groupoid. We \unsure{shall} take note that $\Gamma$ \inserted{does} \uline{not} operate on this groupoid, although it operates on the set of its objects.

\margin{489}